\begin{theorembox}{Theorem (Compact Subsets of $\mathbb{R}$ Are Closed and Bounded)}
\begin{theorem}[Compact Subsets of $\mathbb{R}$ Are Closed and Bounded]
\label{thm:compact-implies-closed-bounded}
Let $K\subseteq\mathbb{R}$.  If $K$ is compact, then $K$ is closed and
bounded.  Explicitly,
\[
  K\text{ compact}
  \Longrightarrow
  \leqft[
  \mathbb{R}\setminus K\text{ is open}
  \right]
  \land
  \leqft[
  (\exists M>0)(\forall x\in K)(|x|\leq M)
  \right].
\]
\hyperref[prf:compact-implies-closed-bounded]{Go to proof.}
\end{theorem}
\end{theorembox}
\begin{remark*}[Standard quantified statement]
\[
K \text{ compact} \Longrightarrow (K \text{ closed} \land K \text{ bounded}).
\]
\end{remark*}
\begin{remark*}[Interpretation]
This is the easy half of Heine--Borel and holds the same way in every metric space. Boundedness comes from covering $K$ by growing balls about a fixed point; closedness comes from separating an outside point from $K$ by disjoint balls. The converse --- closed and bounded implies compact --- is special to $\mathbb{R}^{n}$ and is the next theorem.
\end{remark*}
\begin{dependencies}
\begin{itemize}
  \item \hyperref[def:compact-set]{Compact Set in $\mathbb{R}$}
  \item \hyperref[def:closed-set]{Closed Set in $\mathbb{R}$}
  \item \hyperref[def:bounded-subset-of-r]{Bounded Subset of the Real Line}
\end{itemize}
\end{dependencies}
